\documentclass[11pt,a4paper]{article}
\usepackage[T1]{fontenc}
\usepackage[utf8]{inputenc}
\usepackage{lmodern,amsmath,amssymb,amsthm,booktabs}
\usepackage[margin=27mm]{geometry}
\usepackage[hidelinks]{hyperref}
\hypersetup{pdftitle={A Positive Action Plateau from Classical Velocity Memory},
pdfauthor={navstokgap research working notes}}
\setlength{\emergencystretch}{3em}
\newtheorem{proposition}{Proposition}
\newcommand{\E}{\mathbb E}
\newcommand{\Var}{\operatorname{Var}}
\newcommand{\ha}{\mathsf h}
\title{A Positive Action Plateau from Classical Velocity Memory}
\author{navstokgap research working notes}
\date{7 September 2026}
\begin{document}
\maketitle
\begin{abstract}
A finite reversible classical velocity process determines a positive limiting
action coefficient through its integrated velocity correlation. We define the
coefficient operationally as mass times displacement variance divided by
observation duration, derive its exact spectral formula and bounds, and compute
a two-velocity example. Finite speed forces the short-duration coefficient to
vanish while a positive long-duration plateau survives. This separates
microscopic resolution from observation time and identifies the correlation
law that a universal action-scale proposal must explain.
\end{abstract}

\section{The observable and its classical premises}
The candidate action field is an ensemble observable
\begin{equation}\label{eq:def}
 \ha_\Delta(t)=\frac{m}{\Delta}\Var[X(t+\Delta)-X(t)],\qquad m,\Delta>0.
\end{equation}
Here $X$ is position on the line, $t$ is physical time and $\Delta$ is an
observation duration. Its units are mass times length squared per time, hence
action. For Gaussian increments with variance $\kappa\Delta/m$, this
normalization yields $\ha_\Delta=\kappa$, matching the earlier free-kernel
model. For stationary velocity statistics the observable is independent of
$t$, and we write $\ha(\Delta)$.

The present premises are classical stochastic ones: a finite set of velocities,
ordinary probabilities, finite transition rates, a stationary ensemble and
detailed balance. Their mathematical consequences give a positive transport
scale. The project's stronger target adds a physical origin for these premises,
universality across systems, and a derivation identifying the scale with
$\hbar$. We reserve $h=2\pi\hbar$ for Planck's constants and use $\ha$ for
the candidate observable. An independent dynamical field would require an
additional state variable and its equation of motion.

\section{Finite speed and the correlation formula}
The observation duration determines which part of the velocity dynamics is
resolved. In particular, bounded velocity gives a sharp elementary constraint.

\begin{proposition}[Microscopic limit]\label{prop:short}
If $X$ has absolutely continuous paths with $|\dot X|\le u<c$ almost
everywhere, almost surely, then
\[
 0\le\ha_\Delta(t)\le m u^2\Delta\longrightarrow0
 \quad(\Delta\downarrow0).
\]
\end{proposition}
\begin{proof}
The displacement has absolute value at most $u\Delta$.
Its variance is at most its second moment, which is at most $u^2\Delta^2$.
Divide by $\Delta/m$.
\end{proof}

For a real centered stationary velocity $V$ with finite second moment, set
$X(t)=\int_0^t V(s)\,ds$ and $C(s)=\E[V(s)V(0)]$. Fubini's theorem on
each finite square is justified by Cauchy--Schwarz and stationarity. Symmetry
of the covariance gives
\begin{equation}\label{eq:correlation}
 \ha(\Delta)=2m\int_0^\Delta(1-s/\Delta)C(s)\,ds.
\end{equation}
If $C\in L^1(0,\infty)$, dominated convergence yields
\begin{equation}\label{eq:gk}
 H_*:=\lim_{\Delta\to\infty}\ha(\Delta)
 =2m\int_0^\infty C(s)\,ds=2mD.
\end{equation}
This is the Green--Kubo diffusion coefficient $D$, with an action normalization
\cite[(1.3), Proposition 2.1]{Pavliotis2010}. The finite-state assumptions
below make its strict positivity and convergence explicit.

\section{A positive plateau from finite reversible dynamics}
Let $J_t$ be a stationary irreducible continuous-time Markov chain on
$\{1,\ldots,k\}$, $k\ge2$, with invariant probabilities $\pi_i>0$ and
generator $Q$. Thus $Q_{ij}\ge0$ for $i\ne j$, $Q\mathbf1=0$ and
$\pi_iQ_{ij}=\pi_jQ_{ji}$. All rates are finite. Assign velocities $v_i$ with
\[
 \sum_i\pi_i v_i=0,\qquad \sigma_v^2:=\sum_i\pi_i v_i^2>0,
 \qquad |v_i|\le u<c.
\]
Set $V(t)=v_{J_t}$ and integrate to obtain $X$. The Markov state is $(X,J)$;
the velocity label retains the information required for future motion.
Paths are piecewise linear, with finitely many velocity jumps on compact
intervals almost surely. These jumps represent impulsive interactions rather
than a smooth prescribed force law.

Use $\langle f,g\rangle_\pi=\sum_i\pi_if_ig_i$. On the mean-zero
subspace, $-Q$ is positive definite: detailed balance gives
\[
 \langle f,-Qf\rangle_\pi
 =\tfrac12\sum_{i,j}\pi_iQ_{ij}(f_j-f_i)^2,
\]
and irreducibility makes its nullspace consist of constants. Let its positive
eigenvalues be $\gamma_j$, with an orthonormal eigenbasis $e_j$, and write
$v=\sum_{j=1}^{k-1}a_je_j$.

\begin{proposition}[Classical action plateau]\label{prop:plateau}
Under these premises,
\begin{align}
 \ha(\Delta)&=2m\sum_{j=1}^{k-1}\frac{a_j^2}{\gamma_j}
 \left[1-\frac{1-e^{-\gamma_j\Delta}}{\gamma_j\Delta}\right],
 \label{eq:spectral}\\
 H_*&=2m\langle v,(-Q)^{-1}v\rangle_\pi
 =2m\sum_{j=1}^{k-1}\frac{a_j^2}{\gamma_j}>0.\label{eq:plateau}
\end{align}
The inverse acts on mean-zero functions. Writing
$\gamma_{\min}=\min_j\gamma_j$ and $\gamma_{\max}=\max_j\gamma_j$ gives
\begin{equation}\label{eq:bounds}
 \frac{2m\sigma_v^2}{\gamma_{\max}}\le H_*
 \le\frac{2m\sigma_v^2}{\gamma_{\min}},\qquad
 0\le1-\frac{\ha(\Delta)}{H_*}\le\frac{1}{\gamma_{\min}\Delta}.
\end{equation}
Moreover $\ha(\Delta)$ increases strictly from zero to $H_*$; in particular
$\ha(\Delta)\ge H_*/2$ for $\Delta\ge2/\gamma_{\min}$.
\end{proposition}
\begin{proof}
The transition semigroup gives
$C(s)=\langle v,e^{sQ}v\rangle_\pi=\sum_j a_j^2e^{-\gamma_js}$.
Inserting this into \eqref{eq:correlation} and integrating each exponential
proves \eqref{eq:spectral}. At least one $a_j$ is nonzero, so the limit is
strictly positive. The identity $\sum_j a_j^2=\sigma_v^2$ gives the two
spectral bounds. Each relative correction is bounded above by
$1/(\gamma_{\min}\Delta)$, which gives the convergence estimate.
Finally the scalar function $f(y)=1-(1-e^{-y})/y$ satisfies
$f'(y)=[1-(1+y)e^{-y}]/y^2>0$ for $y>0$, since $e^y>1+y$.
Its endpoint limits are zero and one. Positive weighted summation proves
strict monotonicity and the stated duration-dependent bound.
\end{proof}

This is a finite-state spectral specialization of the generator Poisson
representation in \cite[\S2]{Pavliotis2010}. The bounds concern the observable
\eqref{eq:def} and fixed model parameters. Within the same axiom class,
replacing $Q$ by $rQ$ for $r>0$ changes $H_*$ to $H_*/r$. Thus a uniform
lower bound over models requires control of the transition rates as well as
nonzero velocity variance. The spectral gap of the velocity generator has
inverse-time units; $H_*$ has action units.

\section{Two velocities: the full time dependence}
Choose velocities $\pm u$, $0<u<c$, each with stationary probability $1/2$, and a
reversal rate $\lambda>0$:
\[
 Q=\begin{pmatrix}-\lambda&\lambda\\\lambda&-\lambda\end{pmatrix},
 \qquad V(t)=u\sigma_0(-1)^{N_t}.
\]
Here $N_t$ is a Poisson counting process of rate $\lambda$, independent of
the uniform sign $\sigma_0$. Its generating function gives
$\E[(-1)^{N_s}]=e^{-2\lambda s}$, hence
\begin{equation}\label{eq:telegraph}
 C(s)=u^2e^{-2\lambda s},\qquad
 \ha(\Delta)=\frac{mu^2}{\lambda}
 \left[1-\frac{1-e^{-2\lambda\Delta}}{2\lambda\Delta}\right].
\end{equation}
Therefore
\[
 \ha(\Delta)\sim mu^2\Delta\quad(\Delta\downarrow0),\qquad
 H_*=\frac{mu^2}{\lambda}>0\quad(\Delta\to\infty).
\]
The correlation time is $1/(2\lambda)$, so the limiting action is twice
mass times velocity variance times correlation time. A particle starting at
$X(0)=0$ has support inside $[-u\Delta,u\Delta]$ at time $\Delta$;
the no-reversal events give atoms at both endpoints. The model uses probability
measures, including these atoms, rather than assuming a Gaussian density at
every duration.

The velocity-resolved position measures
$p_\pm(A,t)=\Pr\{X(t)\in A,V(t)=\pm u\}$ satisfy, distributionally,
\[
 \partial_t p_+=-u\partial_xp_+-\lambda p_++\lambda p_-,\qquad
 \partial_t p_-=u\partial_xp_-+\lambda p_+-\lambda p_-.
\]
Adding and subtracting these equations and eliminating their difference gives
the telegraph equation for $p=p_++p_-$,
\[
 \partial_t^2p+2\lambda\partial_tp=u^2\partial_x^2p.
\]
This is the persistent-flight construction associated with the classical
telegraph process \cite[pp.~1--2, (1.1)--(1.3)]{Cinque2022}.

\section{Which limit keeps the action?}
Microscopic resolution and diffusion scaling give different values of the
same observable. Fix $a>0$ with diffusion units and take the formal family
$u_\lambda^2=a\lambda$. For every fixed $\Delta>0$,
\[
 \ha_\lambda(\Delta)=ma
 \left[1-\frac{1-e^{-2\lambda\Delta}}{2\lambda\Delta}\right]
 \longrightarrow ma.
\]
Consequently the iterated limits are
\begin{equation}\label{eq:orders}
 \lim_{\Delta\downarrow0}\lim_{\lambda\to\infty}\ha_\lambda(\Delta)=ma,
 \qquad
 \lim_{\lambda\to\infty}\lim_{\Delta\downarrow0}\ha_\lambda(\Delta)=0.
\end{equation}
This is an exact statement about second moments. The formal scaling sends
$u_\lambda\to\infty$ and therefore leaves any fixed physical speed ceiling.
At fixed $u<c$ and fixed $\lambda$, the positive plateau instead occurs for
observation durations long compared with the correlation time.

\section{Axiom audit and the universality question}
Each premise has a visible mathematical role:
\begin{center}
\begin{tabular}{@{}p{0.36\linewidth}p{0.58\linewidth}@{}}
\toprule
Classical premise & Consequence in the proof\\
\midrule
Stationary centered ensemble & Correlation depends on elapsed time; mean drift is removed\\
Finite irreducible transition dynamics & Inverse generator exists on mean-zero functions\\
Detailed balance & Real orthogonal velocity modes with positive decay rates\\
Nonzero velocity variance & At least one mode contributes to the positive plateau\\
Finite speed & Short-duration action coefficient tends to zero\\
Mass and transition rates & Set the numerical value and units of the plateau\\
\bottomrule
\end{tabular}
\end{center}

These premises establish positivity for each model in the class. Independent
composition permits particles with different transition rates, so universality
adds a further dynamical requirement. For the two-velocity family, a common
limiting value $H_*$ across masses requires
\begin{equation}\label{eq:masslaw}
 \lambda_m=\frac{m u_m^2}{H_*}.
\end{equation}
Equation \eqref{eq:masslaw} states the rate law to explain. It is not supplied
by independent composition. Likewise, the present effective model specifies
impulsive reversals externally; a closed mechanical realization must account
for their momentum exchange and correlation law.

Two reconstruction sources sharpen this requirement. Nelson's stated
Brownian hypothesis supplies $D=\hbar/(2m)$ at the outset
\cite[abstract]{Nelson1966}. Hall and Reginatto instead add momentum
fluctuations whose strength scales inversely with position uncertainty;
their ensemble coefficient $C$ is identified with $(\hbar/2)^2$
\cite[pp.~5, 8]{HallReginatto2002}. These are concrete comparison premises
for the origin of a shared scale.

The next research step is to test whether a mechanical realization can enforce the
mass-weighted correlation integral in \eqref{eq:gk} as a shared invariant.
That would advance the scale-selection target beyond a family of independent
transport coefficients. A separate phase/composition argument must then
identify the common constant with $\hbar$.

\bibliographystyle{plain}
\bibliography{references/library}
\end{document}
